\section{Mô hình và phương pháp}
\label{sec:method}

\subsection{Mô hình hệ thống}
\label{sec:model}

Mạng relay mặt đất sau thiên tai được biểu diễn bằng đồ thị $G=(V,E)$ lấy từ Topology Zoo~\cite{knight2011zoo} và chuẩn hoá vào vùng vuông cạnh $L_{\mathrm{box}}$ với cùng một hệ số tỷ lệ trên hai trục. Nút $c\in V$ là trung tâm cứu hộ, mọi nút còn lại là relay. Cạnh $e\in E$ còn hoạt động có dung lượng backhaul $C_e$; cạnh hỏng có dung lượng bằng $0$. Các điểm phát cảnh báo tạo thành tập $S$; chúng không thuộc $G$ và chỉ truyền được bằng vô tuyến. Mỗi cảnh báo $a\in A$ có nguồn $s_a\in S$, slot phát sinh $r_a$, hạn $d_a$ và kích thước $L_a$ bit. Nhiệm vụ gồm $N$ slot độ dài $\Delta$; mỗi slot chia thành nửa \emph{truy nhập} dài $\tau\Delta$, trong đó nguồn gửi tới relay hoặc UAV, và nửa \emph{downlink} dài $(1-\tau)\Delta$, trong đó UAV gửi xuống một nút của $G$.

UAV bay ở độ cao cố định $H$ theo quỹ đạo $q[0],\dots,q[N]\in\R^2$ thoả
\begin{equation}
\label{eq:mobility}
q[0]=q_s,\qquad q[N]=q_t,\qquad \norm{q[n+1]-q[n]}\le V_{\max}\Delta,\quad n=0,\dots,N-1,
\end{equation}
và vị trí đại diện của slot $n$ là trung điểm $\bar q[n]=\bigl(q[n]+q[n+1]\bigr)/2$.

\paragraph{Kênh.}
Liên kết giữa UAV và điểm mặt đất $k$ ở toạ độ $w_k$ dùng mô hình xác suất tầm nhìn thẳng (PrLoS) của~\cite{huang2025ddatsap}:
\begin{equation}
\label{eq:geometry}
\begin{aligned}
d_k[n]&=\sqrt{\norm{\bar q[n]-w_k}^2+H^2},\qquad
\theta_k[n]=\frac{180}{\pi}\operatorname{atan2}\bigl(H,\norm{\bar q[n]-w_k}\bigr),\\
P_k^{L}[n]&=\frac{1}{1+a\exp\bigl[-b(\theta_k[n]-a)\bigr]},
\end{aligned}
\end{equation}
với độ lợi $g^{L}=\beta_0 d^{-\alpha_L}$ khi có tầm nhìn thẳng và $g^{N}=\mu\beta_0 d^{-\alpha_N}$ khi không. Với công suất phát $p$ và mật độ phổ nhiễu $N_0$, đặt $s=pg/N_0$ (đơn vị Hz); tốc độ khi cấp băng thông $b$ là
\begin{equation}
\label{eq:rate}
R(b,s)=b\log_2\!\Bigl(1+\frac{s}{b}\Bigr),\qquad R(0,s)=0.
\end{equation}
Trong realization $\omega$, trạng thái LoS/NLoS của liên kết tới điểm $k$ ở slot $n$ được lấy mẫu độc lập với xác suất $P_k^{L}[n]$ tính trên chính quỹ đạo đang xét; mọi kế hoạch được so sánh trên cùng các số ngẫu nhiên. Liên kết từ nguồn tới relay mặt đất có độ lợi tất định $\beta_0 d^{-\alpha_G}$ và chỉ tồn tại khi $R(B_{\mathrm{tot}},s)\ge R_{\min}$. Nguồn phát công suất $p_s$; UAV có tổng công suất $P_U$ chia đều cho tối đa hai downlink hoạt động trong một slot, nên ngân sách công suất không tăng theo quy mô.

\paragraph{Đường chuyển tiếp.}
Mỗi cảnh báo có tối đa $K$ candidate path thuộc hai loại: đường \emph{mặt đất} $s_a\to v\leadsto c$ qua một relay $v$ trong tầm, và đường \emph{UAV} $s_a\to\text{UAV}\to u\leadsto c$ với nút thả $u\in V$. Đoạn $v\leadsto c$ là đường backhaul ít chặng nhất trên các cạnh còn hoạt động. Tập candidate gồm tối đa hai đường mặt đất có chi phí ước lượng thấp nhất, bổ sung bằng các đường UAV có chi phí thấp nhất; tập này cố định trước khi lập kế hoạch và mọi phương pháp dùng chung.

\subsection{Phát biểu bài toán}
\label{sec:problem}

Một kế hoạch gồm quỹ đạo $q$ thoả \eqref{eq:mobility}, lựa chọn đường $x_{a,p}\in\set{0,1}$ với $\sum_p x_{a,p}\le 1$ (cho phép không phục vụ cảnh báo), và băng thông $b_{e,n}\ge 0$ trên các liên kết vô tuyến. Kế hoạch được lập trước nhiệm vụ chỉ từ tốc độ kỳ vọng và được đánh giá trên các realization kênh độc lập. Khi thực thi, trong mỗi slot $n$:
\begin{enumerate}
\item \emph{Phổ:} tổng băng thông của các liên kết truy nhập và của các downlink đều không vượt $B_{\mathrm{tot}}$; có tối đa hai downlink hoạt động.
\item \emph{Dung lượng:} liên kết vô tuyến $e$ mang tối đa $t_e\Delta\,R(b_{e,n},s_{e,n})$ bit, với $t_e=\tau$ cho truy nhập và $t_e=1-\tau$ cho downlink; cạnh backhaul mang tối đa $\Delta C_e$ bit, chia chung cho mọi cảnh báo đi qua.
\item \emph{Nhân quả:} bit của $a$ có mặt ở nguồn từ đầu slot $r_a$; truyền trong slot $n$ chỉ dùng dữ liệu có trong buffer ở đầu slot $n$, nên bit nhận trong slot $n$ chỉ được chuyển tiếp từ slot $n+1$; bit chỉ đi trên đường đã chọn và buffer không âm.
\item \emph{Hạn:} ở đầu slot $d_a$, mọi bit còn lại của $a$ bị huỷ.
\end{enumerate}
Mọi phương pháp dùng chung bộ điều phối phục vụ mỗi hàng đợi theo thứ tự hạn sớm nhất (EDF).

\begin{definition}[Cảnh báo đúng hạn]
\label[definition]{def:timely}
Cảnh báo $a$ \emph{đúng hạn} trong realization $\omega$, ký hiệu $z_{a,\omega}=1$, nếu ít nhất $L_a$ bit của nó tới trung tâm $c$ không muộn hơn cuối slot $d_a-1$; ngược lại $z_{a,\omega}=0$.
\end{definition}

Gọi $D_{a,\omega}$ là số bit của $a$ tới $c$ trước hạn. Các kế hoạch được so sánh theo thứ tự từ điển của
\begin{equation}
\label{eq:objective}
\Bigl(\hat P_{\mathrm{timely}},\ \mathrm{Conn},\ -\sum_{\omega}\sum_{a}\frac{L_a-D_{a,\omega}}{L_a}\Bigr),
\qquad
\hat P_{\mathrm{timely}}=\frac{1}{M\abs{A}}\sum_{\omega=1}^{M}\sum_{a\in A}z_{a,\omega},
\end{equation}
trong đó $\mathrm{Conn}$ là tỷ lệ nguồn có đường theo thời gian tới $c$ trên đồ thị mở rộng theo thời gian: mỗi bước truyền tốn một slot, được phép chờ trong buffer, và liên kết UAV dùng được ở slot $n$ khi $R(B_{\mathrm{tot}},s_{e,n})\ge R_{\min}$.

\begin{remark}
Thành phần thứ ba là tỷ lệ dữ liệu còn thiếu tại hạn (shortfall), thay cho tổng độ trễ: vì cảnh báo quá hạn bị huỷ, cảnh báo trễ không có thời điểm hoàn thành để tính độ trễ.
\end{remark}

\subsection{Phương pháp cơ sở}
\label{sec:baselines}

Hai phương pháp cơ sở dùng chung tập candidate, bộ điều phối EDF và quy tắc chia đều $B_{\mathrm{tot}}$ cho các liên kết đang có dữ liệu chờ trong mỗi nửa slot; chúng chỉ khác nhau ở vai trò của UAV. Cả hai lập kế hoạch từ tốc độ kỳ vọng và không dùng các realization đánh giá.

\paragraph{B0 --- không có UAV.}
UAV đứng yên tại điểm xuất phát và không được giao cảnh báo nào. Mỗi cảnh báo đi theo đường mặt đất có chi phí ước lượng thấp nhất; nguồn không có relay nào trong tầm thì cảnh báo không được phục vụ. Với B0, $\mathrm{Conn}$ được tính trên đồ thị không có liên kết UAV. Chi phí lập kế hoạch là $O(\abs{A}K)$.

\paragraph{B1 --- UAV bay theo quỹ đạo cố định.}
UAV xuất phát từ trung tâm, lần lượt bay tới tâm các vùng hư hại theo thứ tự láng giềng gần nhất với tốc độ $V_{\max}$, chia đều thời gian còn lại để lơ lửng tại các vùng, rồi trở về; vùng được thêm sau cùng bị bỏ nếu hành trình không vừa $N$ slot. Với mỗi cảnh báo, mọi candidate được mô phỏng riêng cho cảnh báo đó trên quỹ đạo này với tốc độ kỳ vọng và toàn bộ băng thông; B1 chọn candidate giao xong sớm nhất (hoà thì lấy chỉ số nhỏ hơn) và không phục vụ cảnh báo nếu không candidate nào kịp hạn. Chi phí lập kế hoạch là $O(\abs{A}KNH_{\max})$, với $H_{\max}$ là số chặng lớn nhất của một candidate.

\subsection{Cận trên theo từng realization}
\label{sec:bound}

Cố định quỹ đạo $q$, tập candidate và realization $\omega$; khi đó $s_{e,n}$ đã biết. Với chặng $h$ của candidate $p$ và slot $n\in[r_a+h,\,d_a-1]$, gọi $f_{a,p,h,n}$ là số bit đi qua chặng và $B_{a,p,h,n}$ là buffer ở đầu slot. Chọn các điểm tiếp tuyến $b_k=B_{\mathrm{tot}}2^{-k}$, $k=0,\dots,9$, và xét bài toán quy hoạch tuyến tính
\begin{subequations}
\label{eq:lp}
\begin{align}
\max\quad & \textstyle\sum_{a\in A} z_a \label{eq:lp-objective}\\
\text{s.t.}\quad
& \textstyle\sum_{p} x_{a,p}\le 1, \label{eq:lp-choice}\\
& B_{a,p,0,r_a}=L_a x_{a,p},\qquad B_{a,p,h,r_a+h}=f_{a,p,h-1,r_a+h-1}\ (h\ge 1), \label{eq:lp-init}\\
& B_{a,p,h,n+1}=B_{a,p,h,n}-f_{a,p,h,n}+f_{a,p,h-1,n},\qquad f_{a,p,h,n}\le B_{a,p,h,n}, \label{eq:lp-buffer}\\
& L_a z_a\le \textstyle\sum_{p}\sum_{n} f_{a,p,H_p-1,n}, \label{eq:lp-deadline}\\
& \textstyle\sum_{(a,p,h)\,\text{qua}\,e} f_{a,p,h,n}\le t_e\Delta\bigl(R(b_k,s_{e,n})+R'(b_k,s_{e,n})(b_{e,n}-b_k)\bigr)\quad \forall k, \label{eq:lp-radio}\\
& \textstyle\sum_{(a,p,h)\,\text{qua}\,e} f_{a,p,h,n}\le \Delta C_e, \label{eq:lp-backhaul}\\
& \textstyle\sum_{e\,\text{truy nhập}} b_{e,n}\le B_{\mathrm{tot}},\qquad \sum_{e\,\text{downlink}} b_{e,n}\le B_{\mathrm{tot}}, \label{eq:lp-spectrum}\\
& x,z\in[0,1],\qquad f,B,b\ge 0, \label{eq:lp-domain}
\end{align}
\end{subequations}
trong đó $f_{a,p,-1,n}\equiv 0$, $H_p$ là số chặng của candidate $p$, \eqref{eq:lp-buffer} áp dụng khi $n+1\le d_a-1$, \eqref{eq:lp-radio} áp dụng cho liên kết vô tuyến có $s_{e,n}>0$ (nếu $s_{e,n}=0$ thì luồng qua liên kết bằng $0$) và \eqref{eq:lp-backhaul} cho cạnh backhaul. Gọi $\mathrm{UB}_\omega(q)$ là giá trị tối ưu.

\begin{proposition}[Cận trên]
\label[proposition]{prop:ub}
Với quỹ đạo $q$, tập candidate và realization $\omega$ cố định, $\mathrm{UB}_\omega(q)$ không nhỏ hơn số cảnh báo đúng hạn của mọi kế hoạch có quỹ đạo $q$ khi thực thi trên $\omega$.
\end{proposition}

\begin{proof}
Với $s>0$ cố định, $R(b,s)=b\,\phi(1/b)$ với $\phi(t)=\log_2(1+st)$ lõm, nên $R(\cdot,s)$ là phép phối cảnh của một hàm lõm và do đó lõm trên $b>0$~\cite{boyd2004convex}. Một hàm lõm nằm dưới mọi tiếp tuyến, nên $R(b,s)\le R(b_k,s)+R'(b_k,s)(b-b_k)$ với mọi $k$.
Xét một kế hoạch bất kỳ có quỹ đạo $q$ và quá trình thực thi của nó trên $\omega$. Đặt $x_{a,p}=1$ cho đường được chọn, $z_a=z_{a,\omega}$, $f_{a,p,h,n}$ bằng số bit truyền qua chặng $h$ trong slot $n$, $B_{a,p,h,n}$ bằng buffer ở đầu slot và $b_{e,n}$ bằng băng thông được cấp. Ràng buộc nhân quả và buffer cho \eqref{eq:lp-init}--\eqref{eq:lp-buffer}; \cref{def:timely} cho \eqref{eq:lp-deadline}; ràng buộc phổ và dung lượng cho \eqref{eq:lp-spectrum} và \eqref{eq:lp-backhaul}; bất đẳng thức tiếp tuyến cho \eqref{eq:lp-radio}. Điểm này khả thi và có giá trị mục tiêu bằng số cảnh báo đúng hạn, nên không vượt $\mathrm{UB}_\omega(q)$.
\end{proof}

Bài toán \eqref{eq:lp} còn nới lỏng tính nguyên của $x,z$, thứ tự EDF, giới hạn hai downlink và biết trước kênh; mỗi nới lỏng chỉ làm miền khả thi rộng thêm. Cận chỉ hợp lệ với quỹ đạo và tập candidate đã cho, không phải cận toàn cục trên mọi quỹ đạo. Với mỗi scenario, cận được lấy trung bình trên các realization và chia cho $\abs{A}$; gap của một phương pháp là $(\overline{\mathrm{UB}}-\overline{P})/\overline{\mathrm{UB}}$. B1 được so với cận trên quỹ đạo B1; B0 được so với cận của cùng bài toán khi chỉ giữ các candidate mặt đất.

Để đo độ chặt, trên các instance thu nhỏ bài toán \eqref{eq:lp} được giải lại với $x,z$ nhị phân (MILP). Kế hoạch dựng từ nghiệm MILP --- đường đã chọn và băng thông của nghiệm, giữ hai downlink lớn nhất nếu vượt giới hạn --- được thực thi trên cùng $\omega$, cho chuỗi
\begin{equation}
\label{eq:sandwich}
P_\omega(\text{kế hoạch MILP})\ \le\ \mathrm{OPT}_\omega(q)\ \le\ \mathrm{UB}^{\mathrm{MILP}}_\omega(q)\ \le\ \mathrm{UB}_\omega(q).
\end{equation}
Kế hoạch MILP biết trước kênh nên chỉ dùng để đo độ chặt của cận, không dùng để so sánh phương pháp.

\subsection{Độ khó của bài toán}
\label{sec:hardness}

\begin{proposition}
\label[proposition]{prop:hardness}
Ngay cả khi quỹ đạo cố định, dữ liệu được chia nhỏ tuỳ ý giữa các slot, mạng chỉ có hai relay và mỗi cảnh báo có hai candidate path, bài toán tối đa số cảnh báo đúng hạn là NP-khó.
\end{proposition}

\begin{proof}
Quy dẫn từ PARTITION: cho các số nguyên dương $s_1,\dots,s_m$ có tổng $2Q$, hỏi có tập con có tổng $Q$ hay không. Dựng mạng gồm trung tâm $c$, hai relay $v_1,v_2$ và hai cạnh backhaul $(v_j,c)$ có $\Delta C_{(v_j,c)}=Q$ bit. Một nguồn nằm trong tầm của cả hai relay, với độ lợi đủ lớn để $\tau\Delta\,R(B_{\mathrm{tot}}/2,s)\ge 2Q$. Cảnh báo $i$ có $L_i=s_i$, $r_i=0$, $d_i=2$ và hai candidate: qua $v_1$ hoặc qua $v_2$. Theo quy ước nhân quả, bit của cảnh báo đúng hạn phải được gửi lên relay trong slot $0$ và đi qua backhaul trong slot $1$.
Nếu tồn tại tập con có tổng $Q$, cho các cảnh báo tương ứng đi qua $v_1$, phần còn lại qua $v_2$, và cấp $B_{\mathrm{tot}}/2$ cho mỗi liên kết truy nhập trong slot $0$: mỗi liên kết truy nhập mang được $2Q$ bit và mỗi cạnh backhaul mang đúng $Q$ bit trong slot $1$, nên cả $m$ cảnh báo đúng hạn.
Ngược lại, nếu cả $m$ cảnh báo đúng hạn thì tổng kích thước các cảnh báo qua $v_j$ không vượt $Q$ với $j=1,2$, trong khi tổng hai phần bằng $2Q$; vậy mỗi phần bằng $Q$ và PARTITION có lời giải. Phép dựng thực hiện trong thời gian đa thức, nên quyết định ``có ít nhất $m$ cảnh báo đúng hạn hay không'' là NP-khó.
\end{proof}

\begin{remark}
\Cref{prop:hardness} chỉ khẳng định NP-khó theo nghĩa yếu; báo cáo không khẳng định bài toán thuộc NP (băng thông là biến liên tục) và không xét trường hợp $K=1$. Độ khó đến từ ràng buộc mỗi cảnh báo đi theo đúng một đường, không đến từ việc lập lịch trên một liên kết: khi chỉ có một liên kết và dữ liệu được chia nhỏ theo slot, bài toán tương ứng với lập lịch một máy có ngắt quãng $1\,|\,r_j,\mathrm{pmtn}\,|\,\sum U_j$, vốn giải được trong thời gian đa thức~\cite{lawler1990preemptive}. Phiên bản không ngắt quãng $1\,|\,r_j\,|\,\sum U_j$ là NP-khó mạnh~\cite{lenstra2020elements}, nhưng mô hình ở đây cho phép chia bit theo slot nên kết quả đó không áp dụng trực tiếp.
\end{remark}

\subsection{Độ phức tạp tính toán}
\label{sec:complexity}

Sinh candidate path tốn $O(\abs{V}+\abs{E})$ cho một lần tìm kiếm theo chiều rộng từ $c$ và $O(\abs{A}\abs{V}\log\abs{V})$ cho việc xếp hạng. Đánh giá một kế hoạch trên $M$ realization tốn $O\bigl(MN(\abs{\mathcal L}+\abs{A}\log\abs{A})\bigr)$, với $\mathcal L$ là tập liên kết; bộ kiểm tra độc lập có cùng bậc, còn $\mathrm{Conn}$ tốn $O\bigl(MN(\abs{V}+\abs{E}+\abs{S}\abs{V})\bigr)$. B0 lập kế hoạch trong $O(\abs{A}K)$ và B1 trong $O(\abs{A}KNH_{\max})$. Bài toán \eqref{eq:lp} có $\sum_a\sum_p H_p(d_a-r_a)$ biến luồng và cùng số biến buffer, tối đa $\abs{\mathcal L_{\mathrm{vt}}}N$ biến băng thông và tối đa $10\abs{\mathcal L_{\mathrm{vt}}}N$ ràng buộc tiếp tuyến, với $\mathcal L_{\mathrm{vt}}$ là các liên kết vô tuyến của tập candidate. Cận được giải bằng HiGHS qua \texttt{scipy}; báo cáo không đưa ra cận thời gian đa thức cho bộ giải mà đo thời gian thực của mọi bài toán LP (\cref{sec:experiments}).

\subsection{Phương pháp tìm kiếm theo khối và sửa lịch}
\label{sec:algorithm}

\paragraph{Vì sao các khối không là bài toán lồi.}
Cấu trúc ba khối của~\cite{huang2025ddatsap} được giữ lại, nhưng mỗi khối không được giải bằng một bài toán lồi. Ba thử nghiệm sơ bộ trên sáu scenario đánh giá với mười realization cho thấy lý do. Băng thông lấy từ nghiệm LP ở \cref{sec:bound}, dùng làm lịch tĩnh, làm tỷ lệ đúng hạn giảm từ \num{0.410} của B1 xuống \num{0.128}; dùng làm trọng số chia băng thông đạt \num{0.404}; đường chọn theo nghiệm LP đạt \num{0.412}. Như vậy chỉ dẫn từ LP không chuyển được sang bộ điều phối EDF, còn khối quỹ đạo dưới kênh PrLoS không có hàm thay thế lồi hợp lệ cho mục tiêu đếm cảnh báo. Ngược lại, tìm kiếm cục bộ được chấm điểm bằng chính bộ đánh giá nâng tỷ lệ này lên \num{0.427} khi chấm trên tốc độ kỳ vọng và \num{0.449} khi chấm trên năm realization lấy mẫu; chấm trên tốc độ kỳ vọng còn làm một scenario giảm từ \num{0.372} xuống \num{0.305}. Vì vậy mỗi khối ở đây là một bước tìm kiếm cục bộ có đánh giá.

\paragraph{Chia băng thông theo trọng số.}
Kế hoạch mang trọng số $w_{e,n}\ge 0$ cho mỗi liên kết vô tuyến $e$ và slot $n$. Trong mỗi nửa slot, gọi $\mathcal B_n$ là tập liên kết truy nhập có dữ liệu chờ, hoặc tập tối đa hai downlink có hạn đầu hàng đợi sớm nhất; băng thông được cấp là
\begin{equation}
\label{eq:weights}
b_{e,n}=B_{\mathrm{tot}}\frac{w_{e,n}}{\sum_{e'\in\mathcal B_n}w_{e',n}},\qquad e\in\mathcal B_n,
\end{equation}
\noindent và chia đều khi tổng trọng số bằng $0$. Quy tắc chia đều của B0, B1 là trường hợp mọi trọng số bằng $1$. Trọng số cố định trước nhiệm vụ, còn băng thông thực tế chỉ phụ thuộc hàng đợi hiện tại, nên chính sách không dùng thông tin tương lai.

\paragraph{Chấm điểm thiết kế.}
Phương pháp tìm kiếm không bao giờ dùng các realization đánh giá. Nó chấm một kế hoạch trên $M_d=5$ realization thiết kế, lấy từ một luồng số ngẫu nhiên khác với luồng đánh giá, theo khoá từ điển \eqref{eq:objective} cộng trên các realization đó. Mỗi lần chấm là một lần đánh giá, và tổng số lần đánh giá bị giới hạn bởi ngân sách $E=2000$; khi hết ngân sách, phương pháp trả về kế hoạch tốt nhất đang có. Một thay đổi chỉ được nhận khi khoá tăng chặt, nên khoá của kế hoạch hiện tại không bao giờ giảm. Khi chấm thiết kế, bộ kiểm tra độc lập được bỏ qua và $\mathrm{Conn}$ được ghi nhớ theo quỹ đạo và realization, vì cả hai không ảnh hưởng tới quyết định tìm kiếm; cách này giữ nguyên kế hoạch trả về và nhanh hơn khoảng \num{2.9} lần. Lần đánh giá cuối trên các realization đánh giá vẫn chạy bộ kiểm tra.

\paragraph{Ba khối của B2.}
B2 khởi tạo từ kế hoạch B1 với mọi trọng số bằng $1$ và lặp tối đa mười vòng. \emph{Khối đường} duyệt cảnh báo theo hạn tăng dần và thử lần lượt các lựa chọn không phục vụ và từng candidate. \emph{Khối băng thông} duyệt các liên kết vô tuyến đang được dùng theo tổng kích thước cảnh báo đi qua, giảm dần, và thử nhân toàn bộ trọng số của liên kết với $2$, rồi với $1/2$ nếu phép nhân đôi bị từ chối. \emph{Khối quỹ đạo} duyệt một họ hành trình khả thi theo cấu tạo: UAV rời trung tâm, bay thẳng với tốc độ $V_{\max}$ qua một tập con có thứ tự của các vùng hư hại, lơ lửng tại mỗi điểm dừng rồi quay về. Thời gian còn dư được chia cho các điểm dừng đều nhau, theo số cảnh báo của vùng, hoặc theo tổng $1/(d_a-r_a)$ của vùng; điểm dừng là tâm vùng hoặc trọng tâm các nguồn trong vùng. Với ba vùng, họ này có tối đa $90$ hành trình và chứa quỹ đạo của B1. Vòng lặp dừng khi một vòng không nhận thay đổi nào.

\paragraph{B3: mục tiêu leximin.}
B3 giống B2 nhưng dùng khoá khác. Với $\bar f_a$ là tỷ lệ bit của $a$ tới trung tâm trước hạn, lấy trung bình trên các realization thiết kế, tốc độ giao trung bình của $a$ chia cho tốc độ yêu cầu $L_a/((d_a-r_a)\Delta)$ bằng đúng $\bar f_a$, nên đây là phiên bản chuẩn hoá của mục tiêu max--min của~\cite{huang2025ddatsap}. Vì chỉ một cảnh báo không thể giao cũng làm $\min_a\bar f_a=0$ với mọi kế hoạch, B3 so sánh theo leximin:
\begin{equation}
\label{eq:leximin}
\bigl(\bar f_{(1)},\bar f_{(2)},\dots,\bar f_{(\abs{A})},\ \textstyle\sum_{\omega}\sum_a z_{a,\omega},\ \mathrm{Conn}\bigr),\qquad \bar f_{(1)}\le\bar f_{(2)}\le\dots\le\bar f_{(\abs{A})},
\end{equation}
\noindent trong đó cảnh báo không được phục vụ có $\bar f_a=0$. Chênh lệch giữa B2 và B3 đo cái giá của việc tối ưu công bằng thay cho số cảnh báo đúng hạn.

\paragraph{P: sửa lịch theo nguy cơ trễ hạn.}
P là B2 thêm một vòng sửa lịch sau khối quỹ đạo của mỗi vòng lặp. Một lần đánh giá có ghi sổ trên kế hoạch hiện tại cho biết, với mỗi cảnh báo và realization thiết kế, độ dư thời gian $d_a-1-t_a$ (bằng $-1$ nếu trễ, với $t_a$ là slot giao xong), chặng giữ dữ liệu của cảnh báo trong nhiều slot nhất và các liên kết đã dùng hết dung lượng. Nguy cơ của một cảnh báo là tỷ lệ realization thiết kế trong đó nó trễ hoặc có độ dư không quá $\theta=2$ slot; tối đa $R=10$ cảnh báo có nguy cơ cao nhất được xét. \emph{Thao tác 1} nhân đôi trọng số của chặng thắt cổ chai, nếu đó là liên kết vô tuyến, chỉ trong các slot $[r_a,d_a)$, nên băng thông được chuyển từ các liên kết khác cùng nửa slot sang cảnh báo này. \emph{Thao tác 2} thử các candidate khác của cảnh báo theo số liên kết nghẽn trong cửa sổ của nó, rồi theo chi phí ước lượng, và nhận candidate đầu tiên làm khoá tăng. Mọi thao tác dùng chung ngân sách với ba khối, nên P và B2 có cùng số lần đánh giá tối đa. \Cref{alg:bcd-repair} tóm tắt phương pháp.

\begin{algorithm}[htbp]
\caption{Tìm kiếm theo khối có đánh giá, kết hợp sửa lịch (B2 khi tắt sửa lịch, P khi bật)}
\label{alg:bcd-repair}
\begin{algorithmic}[1]
\Require Scenario, tập candidate, realization thiết kế, ngân sách $E$, số vòng tối đa $I_{\max}=10$
\Ensure Quỹ đạo $q$, lựa chọn đường $x$, trọng số băng thông $w$
\State $(q,x)\gets$ kế hoạch B1; $w\gets 1$; $\kappa\gets$ khoá của $(q,x,w)$
\For{$i=1,\dots,I_{\max}$}
    \State Khối đường: với từng cảnh báo theo hạn, thử không phục vụ và từng candidate
    \State Khối băng thông: với từng liên kết đang dùng, thử $w_e\gets 2w_e$, sau đó $w_e\gets w_e/2$
    \State Khối quỹ đạo: thử từng hành trình trong họ hành trình khả thi
    \If{sửa lịch được bật}
        \State Đánh giá có ghi sổ; chọn tối đa $R$ cảnh báo có nguy cơ trễ hạn cao nhất
        \ForAll{cảnh báo $a$ có nguy cơ}
            \State Thao tác 1: nhân đôi trọng số chặng thắt cổ chai trong $[r_a,d_a)$
            \State Thao tác 2: thử các candidate khác theo số liên kết nghẽn, rồi chi phí
        \EndFor
    \EndIf
    \State Mỗi phép thử được nhận khi khoá tăng chặt; dừng ngay khi hết ngân sách $E$
    \If{vòng này không nhận thay đổi nào} \textbf{break} \EndIf
\EndFor
\State \Return $(q,x,w)$
\end{algorithmic}
\end{algorithm}

\paragraph{Độ phức tạp.}
Một vòng lặp dùng tối đa $(K+1)\abs{A}+2\abs{\mathcal L_{\mathrm{used}}}+90$ lần đánh giá cho ba khối, cộng $1+RK$ lần cho vòng sửa lịch của P, với $\mathcal L_{\mathrm{used}}$ là các liên kết vô tuyến đang được dùng. Mỗi lần đánh giá tốn $O\bigl(M_dN(\abs{\mathcal L}+\abs{A}\log\abs{A})\bigr)$, còn phát hiện nguy cơ duyệt sổ ghi trong $O\bigl(M_d(N\abs{\mathcal L}+T)\bigr)$ với $T$ là số bản ghi truyền. Ngân sách chặn tổng số lần đánh giá, nên thời gian lập kế hoạch của B2, B3 và P là $O\bigl(EM_dN(\abs{\mathcal L}+\abs{A}\log\abs{A})\bigr)$, không phụ thuộc số vòng lặp thực tế. Phương pháp chỉ bảo đảm khoá của kế hoạch không giảm; nó không bảo đảm tối ưu, và cận LP ở \cref{sec:bound}, giải trên chính quỹ đạo mà mỗi phương pháp chọn, đo phần còn có thể cải thiện.
